\section{2024-11-25 Herrausforderungen Deutschland}

\startitemize[packed]
\item China
\item USA
\item Russland - Ukraine
\item Israel / Palestina
\item EU
\stopitemize

\subsection{S. 120 M1}

{\bf Neue Härte in der Weltpolitik}

\margintext{
{\bf Autoritarismus}

gilt in der Politikwissenschaft als eine diktatorische Form der Herrschaft, welche zwischen Demokratie und dem diktatorischen Totalitarismus liegt.
}

\startitemize[packed]
\item China {\bf nur} als Ort für deutsche Geschäfte
\startitemize[packed]
\item Autoritatismus wurde vernachlässigt
\item Hohe abhängigkeit von chinesischen Markt bedroht deutschen Wohlstand
\stopitemize
\item Russland
\startitemize[packed]
\item 2014 als Wendepunkt
\item Russland veränderte gewaltsam Grenzen
\item wir zu naive
\stopitemize
\item Migration aus gescheiterten Staaten können die gesellschaftliche Stabilität gefährden (Z. 46-47)
\item Deutschland ist gut vernetzt
\stopitemize

\subsection{S. 121 M2}

\subsection{EU}

\startitemize[packed]
\item Staaten handeln nationalistischer
\item Keine zusammenarbeit
\item Einflusverlusst in einer multipolaren Welt
\item Viele Unstümmigkeiten
\stopitemize

\subsection{China}

\startitemize[packed]
\item  Beziehung zu Deutschland
\startitemize[packed]
\item Großer erxport Partner für Deutschland
\stopitemize
\item Naiver Gedanke, China wird selber gezwungenermaßen Demokratisch
\item Was passiert, wenn China Taiwan angreift
\stopitemize
